\documentclass[11pt]{article}
\usepackage[margin=1in]{geometry}
\usepackage{amsmath,amssymb,amsthm,mathtools}
\usepackage[T1]{fontenc}
\usepackage{lmodern}
\usepackage[colorlinks=true,linkcolor=blue,citecolor=blue,urlcolor=blue]{hyperref}
\newtheorem{theorem}{Theorem}
\newtheorem{lemma}[theorem]{Lemma}
\newcommand{\R}{\mathbb R}
\newcommand{\Om}{\Omega}
\newcommand{\F}{\mathcal F}
\title{Two positive Fisher--KPP equilibria on a uniformly smooth unbounded domain}
\author{Kaiwen Jin\thanks{Email: \href{mailto:yui.ui.siki@gmail.com}{yui.ui.siki@gmail.com}.}}
\date{15 October 2026}
\begin{document}
\maketitle

\begin{abstract}
We construct a connected unbounded domain in $\R^{21}$ with uniformly
$C^{2,1/2}$ boundary on which the Dirichlet Fisher--KPP equation has two
distinct bounded positive equilibria. The boundary consists of two smooth
global graphs whose separation tends to the critical slab width.
A corrected transverse mode gives a slowly decaying subsolution, while a
compact enlargement of the domain permits a second equilibrium with a
faster upper bound. The example gives a negative answer to the unrestricted
strong-KPP uniqueness question. No optimal dimension is asserted.
\end{abstract}

\section{The uniqueness question}

Positive stationary populations with diffusion and density-dependent growth
are described by
\begin{equation}\label{eq:problem}
 -\Delta u=f(u)\quad\hbox{in }\Om,\qquad u=0\quad\hbox{on }\partial\Om.
\end{equation}
Here $\Om$ is the habitat and the boundary is absorbing. A decreasing
per-capita growth rate, meaning that $s\mapsto f(s)/s$ is strictly decreasing
on $(0,1]$, is called the strong-KPP condition here. It gives uniqueness of a positive
equilibrium on smooth bounded domains. On unbounded domains, positive
solutions may approach zero at infinity, and a global comparison between
two solutions is harder to obtain.

Berestycki and Graham proved uniqueness on general uniformly smooth domains
under a nondegeneracy condition on the principal limit spectrum
\cite[Theorem~1.1]{BGKPP}, and asked whether the condition could be removed
\cite[Conjecture~1.3]{BGKPP}. Their later paper retains the general question
\cite[Section~8.1, Open Question~3]{BGSC}. Recently, Rossi and Tough proved uniqueness for the quadratic Fisher--KPP
reaction through dimension 6, and constructed homogeneous counterexamples
on locally analytic domains in every dimension at least 7
\cite[Theorems~1.5, 1.8 and 7.1]{RT}. Their homogeneous construction uses
spherical obstacles with radii tending to zero. The example below instead
has two global boundary graphs with uniform derivative bounds.

To keep both the reaction and its derivative bounded on $[0,\infty)$
while retaining the logistic reaction on $[0,1]$, set
\begin{equation}\label{eq:reaction}
 f(s)=\begin{cases}s-s^2,&0\leq s\leq1,\\s^{-1}-1,&s\geq1.\end{cases}
\end{equation}
This function is bounded with bounded Lipschitz derivative. In particular,
$f\in C^{1,1/2}([0,\infty))$, $f(0)=f(1)=0$, $f'(0)=1$, $f>0$ on $(0,1)$,
$f<0$ on $(1,\infty)$, and $f(s)/s=1-s$ strictly decreases on $(0,1]$.

\begin{theorem}\label{thm:main}
There is a connected unbounded domain $\Om\subset\R^{21}$ with uniformly
$C^{2,1/2}$ boundary for which \eqref{eq:problem}, with \eqref{eq:reaction},
has two distinct solutions $U,T\in C^{2,1/2}_{\mathrm{loc}}(\overline\Om)$
satisfying $0<U,T\leq1$ in $\Om$.
The same functions solve $-\Delta u=u-u^2$ with zero Dirichlet data.
\end{theorem}

The Dirichlet condition is the admissible value $\beta=0$ of
$\beta\partial_nu+(1-\beta)u=0$. Thus one example in Theorem~\ref{thm:main}
refutes the universal uniqueness assertion for all allowed dimensions,
domains, reactions and boundary parameters. It makes no classification
claim for other dimensions or for Robin boundary conditions.

The domain is a strip whose width approaches the critical value $\pi$ with
an inverse-square excess, together with a compact central enlargement.
We first construct an equilibrium bounded below by a multiple of
$|x|^{-2}$. The enlargement also supports a compact positive subsolution.
An exterior supersolution of order $|x|^{-4}$ can be joined to the first
equilibrium and used to construct the second one. The different bounds at
infinity separate the two equilibria.

\section{Comparison functions on an unbounded strip}

We use the polynomial residual
$\F(v)=\Delta v+v-v^2$. A continuous, locally $H^1$ function is a subsolution
if $\F(v)\geq0$ in distributions, and a supersolution if $\F(v)\leq0$.
All subsolutions and supersolutions used to construct equilibria lie in
$[0,1]$, where the reaction in \eqref{eq:reaction} agrees with the polynomial.
The following existence fact records how the boundary at infinity is treated.

\begin{lemma}\label{lem:existence}
Let $\ell:\R^{20}\to(0,\infty)$ be smooth, bounded above and below by positive
constants, with bounded derivatives, and put
$D=\{(x,y):|y|<\ell(x)\}$.
Suppose $a,c$ are continuous and locally $H^1$ up to $\partial D$,
$0\leq a\leq c\leq1$, $\F(a)\geq0$, $\F(c)\leq0$, and $a=0$ on $\partial D$.
Then there is a classical solution $u$ of
$\Delta u+u-u^2=0$ on $D$, with $a\leq u\leq c$ and $u=0$ on $\partial D$.
It is locally $C^{2,1/2}$ up to the boundary. If $a\not\equiv0$, then $u>0$.
\end{lemma}

\begin{proof}
Let $D_n=D\cap\{|x|<n\}$. This is a bounded Lipschitz domain: its end face
meets the two graph faces transversely. Prescribe the boundary extension $a$,
so that the trace is zero on the true boundary and $a$ on the end face.
Set $u_0=a$ and solve
\begin{equation}\label{eq:iteration}
 (-\Delta+1)u_{j+1}=G(u_j),\qquad G(t)=2t-t^2,
 \qquad u_{j+1}-a\in H^1_0(D_n).
\end{equation}
The weak linear Dirichlet theorem \cite[Theorem~8.9]{GT} applies with
constant principal coefficients and zero-order coefficient $-1$ after
reversing signs. Testing a difference with its positive part proves the
comparison principle for $-\Delta+1$. Since $G$ increases on $[0,1]$,
the subsolution and supersolution inequalities give
$a=u_0\leq u_1\leq\cdots\leq c$.

For fixed $n$, let $v$ be the almost everywhere monotone limit.
The bounded functions $G(u_j)$ converge to $G(v)$ in $L^2(D_n)$.
Let $u_*$ solve the linear problem with this limiting right hand side and
the same boundary extension. The energy identity for the difference gives
\[
 \|u_{j+1}-u_*\|_{H^1(D_n)}
 \leq\|G(u_j)-G(v)\|_{L^2(D_n)}\longrightarrow0.
\]
Comparison with $a,c$ identifies $u_*=v$ and gives a weak semilinear solution
$u_n$ between the two functions.

On every fixed interior patch, or fixed true-boundary patch away from the
end face, $0\leq u_n\leq1$ and
$\Delta u_n=-u_n+u_n^2$ has bounded right hand side.
Local weak boundary regularity \cite[Corollary~8.36]{GT}, with zero boundary
data and constant coefficients, gives uniform local $C^{1,1/2}$ bounds.
The right hand side is therefore uniformly $C^{0,1/2}$ on smaller patches.
Theorem~8.9 also gives local $H^2$ regularity; Theorem~9.19, with $k=1$,
then gives interior $C^{2,1/2}$ regularity. Lemma~6.18 and the local estimates
of Corollaries~6.3 and~6.7 in \cite{GT} give uniform $C^{2,1/2}$ bounds on
smaller true-boundary and interior patches. All these constants concern fixed
patches; the distant end face does not enter them.

A diagonal subsequence converges locally in $C^2$, including compact portions
of the true boundary. Its limit satisfies the equation, ordering and zero
boundary condition. If $a$ is nonzero, the limit is nonzero. Since
$\Delta u=-u(1-u)\leq0$, the strong minimum principle on the connected domain
gives $u>0$. No global energy norm or prescribed decay at infinity is used.
\end{proof}

Two elementary facts will also be useful. If a smooth subsolution $v$ vanishes
on the boundary of a subdomain $E$ and $\partial_n v\leq0$, then its zero
extension $\widetilde v$ is a weak subsolution. Indeed, for every nonnegative
compactly supported smooth test function $\zeta$,
\begin{equation}\label{eq:extension}
 \int\bigl(\widetilde v\Delta\zeta+
              (\widetilde v-\widetilde v^2)\zeta\bigr)
 =\int_E\F(v)\zeta-\int_{\partial E}\partial_n v\,\zeta\geq0.
\end{equation}
The formula holds face by face for a cube. The extensions are locally $H^1$
because their traces agree with zero.

The minimum of two smooth supersolutions is a weak supersolution. To see
this without a monotonicity assumption on the reaction, approximate the
minimum of $v,w$ by
$m_\varepsilon=(v+w-\sqrt{(v-w)^2+\varepsilon^2})/2$.
Differentiation gives
$\Delta m_\varepsilon\leq\theta_\varepsilon\Delta v+
(1-\theta_\varepsilon)\Delta w$, where $0\leq\theta_\varepsilon\leq1$.
Thus $\F(m_\varepsilon)$ is bounded above by
\[
 (m_\varepsilon-m_\varepsilon^2)
 -\theta_\varepsilon(v-v^2)-(1-\theta_\varepsilon)(w-w^2).
\]
This tends to zero pointwise and is bounded on every compact set, including
where $v=w$. Testing and passing to the limit proves the assertion.
The minimum is locally $H^1$ by the weak chain rule
\cite[Section~7.4 and pp.~153--154]{GT}.

\section{Two different rates of decay}

Write $x\in\R^{20}$, $y\in\R$, and fix
\begin{equation}\label{eq:geometry}
 R=10^6,\quad b=20000,\quad q=R^2+|x|^2,\quad
 k=1-\frac{24}{q},\quad \ell_0=\frac{\pi}{2k},\quad
 \phi(x,y)=\cos(k(x)y).
\end{equation}
Let $\Om_0=\{(x,y):|y|<\ell_0(x)\}$. Then $q\geq10^{12}$,
$1/2\leq k<1$, $|y|<\pi$ and $0<\phi\leq1$ in $\Om_0$.
The transverse mode $\phi$ vanishes on its boundary. We need to retain the
derivatives of this mode in $x$, since their signs matter near the boundary.

\begin{lemma}\label{lem:barriers}
On $\Om_0$, the functions
\[
 S=q^{-1}\bigl(\phi+bq^{-2}\phi^2\bigr),\qquad
 V_B=Bq^{-2}\phi\quad(B>0)
\]
satisfy $0<S\leq1$, $\F(S)\geq0$ and $\F(V_B)\leq0$.
They vanish on $\partial\Om_0$, and $\partial_nS<0$ there.
\end{lemma}

\begin{proof}
For a positive amplitude $H=H(x)$, direct differentiation gives
\begin{align}\label{eq:residual}
 \frac{\F(H\phi)}{H}
 &=\phi\left(\frac{\Delta_xH}{H}+1-k^2-y^2|\nabla k|^2-H\phi\right)\notag\\
 &\quad-y\sin(ky)\left(\Delta_xk+2\nabla\log H\cdot\nabla k\right).
\end{align}
For $H=Aq^{-p/2}$ and $r=|x|$,
\begin{align}\label{eq:derivatives}
 \frac{\Delta_xH}{H}&=-\frac{20p}{q}+\frac{p(p+2)r^2}{q^2},\notag\\
 \nabla k&=\frac{48x}{q^2},\qquad
 \Delta_xk=\frac{960}{q^2}-\frac{192r^2}{q^3},\notag\\
 \Delta_xk+2\nabla\log H\cdot\nabla k
 &=\frac{960}{q^2}-\frac{96(p+2)r^2}{q^3}.
\end{align}
For $p=4$, the first linear coefficient in \eqref{eq:residual} is at most
$(-56+48)/q=-8/q$, while the last bracket is at least $384/q^2$.
Both $\phi$ and $y\sin(ky)$ are nonnegative. The other terms have nonpositive
sign, proving $\F(V_B)\leq0$ for every amplitude $B>0$.

For the slow subsolution put $h=q^{-1}$ and $\delta=bq^{-2}$.
The last bracket in \eqref{eq:residual} lies between $576/q^2$ and $960/q^2$.
Using $\pi<4$ and $|\nabla k|^2\leq2304q^{-3}$ gives
\begin{equation}\label{eq:uncorrected}
 \F(h\phi)\geq h\phi\left(\frac7q-\frac{576}{q^2}-\frac{36864}{q^3}\right)
                   -\frac{3840h}{q^2}
 \geq\frac{6h\phi}{q}-\frac{3840h}{q^2}.
\end{equation}
The negative term need not contain $\phi$, which is why the squared mode is
added. With $g=\delta h=bq^{-3}$, its exact contribution is
\begin{equation}\label{eq:correction}
 \F(S)=\F(h\phi)+g\bigl[2k^2+(1-4k^2)\phi^2\bigr]+E,
 \quad E=\Delta_x(g\phi^2)-2\delta h^2\phi^3-\delta^2h^2\phi^4.
\end{equation}
To bound the remainder uniformly, \eqref{eq:derivatives} gives
\[
 |\nabla_x\phi|\leq192q^{-3/2},\quad
 |\Delta_x\phi|\leq40704q^{-2},\quad
 |\nabla_x(\phi^2)|\leq384q^{-3/2},\quad
 |\Delta_x(\phi^2)|\leq155136q^{-2}.
\]
Also $|\nabla g|\leq6bq^{-7/2}$ and $|\Delta_xg|\leq168bq^{-4}$.
The product rule and the two nonlinear terms in \eqref{eq:correction} yield
\begin{equation}\label{eq:error}
 |E|\leq(159914b+b^2)q^{-4}<4\cdot10^9\frac{h}{q^3}.
\end{equation}
For $\phi\leq1/4$, the square bracket in \eqref{eq:correction} is at least
$5/16$, and hence
\[
 \F(S)\geq\left(6250-3840-\frac{4\cdot10^9}{q}\right)\frac{h}{q^2}>0.
\]
For $\phi\geq1/4$ the bracket is at least $-3$, giving
\[
 \F(S)\geq\left(\frac32-\frac{63840}{q}-\frac{4\cdot10^9}{q^2}\right)
                 \frac hq>0.
\]
Both inequalities hold for all $q\geq10^{12}$.
Finally $S\leq q^{-1}(1+bq^{-2})<1$. On the boundary its normal derivative
is $h\partial_n\phi=-h|\nabla\phi|<0$, since
$|\partial_y\phi|=k$ there.
\end{proof}

\section{A compact enlargement and two equilibria}

Choose a smooth radial cutoff $0\leq\chi\leq1$ which is 1 for $r\leq2000$
and 0 for $r\geq4000$, and set
\begin{equation}\label{eq:domain}
 \ell(x)=\ell_0(x)+2000\chi(x),\qquad
 \Om=\{(x,y):|y|<\ell(x)\}.
\end{equation}
For example, if $\eta(t)=e^{-1/t}$ for $t>0$ and $\eta(t)=0$ otherwise, one
can take
$\chi=\eta(4000^2-r^2)/[\eta(4000^2-r^2)+\eta(r^2-2000^2)]$.
Its denominator is positive. The two boundary graphs have bounded derivatives
through order 3 and have separation at least $\pi$. They therefore admit
common-radius $C^{2,1/2}$ charts with uniformly bounded norms. The domain is
connected through its central plane $y=0$, is unbounded, and agrees with
$\Om_0$ for $r\geq4000$.

Extend $S$ by zero from $\Om_0$ into $\Om$. By
\eqref{eq:extension} and Lemma~\ref{lem:barriers}, it is a weak subsolution.
It lies below the constant supersolution 1. Lemma~\ref{lem:existence} gives
an equilibrium $U$ on $\Om$ with
\begin{equation}\label{eq:slow}
 0<U\leq1,\qquad U(x,0)\geq q^{-1}.
\end{equation}

We next construct a supersolution whose exterior bound is faster.
On the compact collar
$4500\leq r\leq5500$, $|y|\leq\ell_0(x)$, let
$M=\sup|\partial_yU|<\infty$.
For $d=\ell_0(x)-|y|$, integration from the nearest endpoint gives $U\leq Md$.
On the other hand,
\[
 \phi=\sin(kd)\geq\frac{2kd}{\pi}\geq\frac d\pi.
\]
The finite choice
$B=1+\pi M(R^2+5500^2)^2$ consequently gives $V_B\geq U$ on this collar.
Define
\begin{equation}\label{eq:patched}
 W=\begin{cases}U,&r\leq5000,\\ \min(U,V_B),&r>5000.\end{cases}
\end{equation}
It equals $U$ on a neighborhood of $r=5000$. The minimum property proved
above and Lemma~\ref{lem:barriers} show that $W$ is a weak supersolution on
the whole of $\Om$. Moreover $0<W\leq1$, it vanishes on the true boundary,
and
\begin{equation}\label{eq:fast}
 W\leq Bq^{-2}\phi\qquad(r>5000).
\end{equation}

To obtain a positive solution below $W$, consider the cube $Q=(-100,100)^{21}$.
Its closure lies strictly inside $\Om$: on it $r<2000$ and $|y|\leq100$,
whereas $\ell\geq2000+\pi/2$.
The function
\[
 \psi(z)=\prod_{i=1}^{21}\cos(\pi z_i/200)\quad(z\in Q)
\]
satisfies $0<\psi\leq1$, zero boundary data and
$-\Delta\psi=\lambda_Q\psi$, with $\lambda_Q=21\pi^2/40000<1$.
Its outward derivative is nonpositive on every face.
Let $m_Q=\min_{\overline Q}W>0$ and
$\varepsilon=\min((1-\lambda_Q)/2,m_Q/2)>0$.
The zero extension of $\varepsilon\psi$ lies below $W$, and in $Q$
\[
 \F(\varepsilon\psi)=\varepsilon\psi
                   \bigl[(1-\lambda_Q)-\varepsilon\psi\bigr]\geq0.
\]
Formula~\eqref{eq:extension}, applied to the faces, proves its global weak
subsolution inequality. Lemma~\ref{lem:existence} now gives a positive
equilibrium $T$ with $\varepsilon\psi\leq T\leq W\leq1$ and zero Dirichlet data.
For $r>5000$ and $q>B$, \eqref{eq:slow} and \eqref{eq:fast} give
\[
 U(x,0)\geq q^{-1}>Bq^{-2}\geq T(x,0).
\]
Such points exist because $B$ is fixed and finite. The two equilibria are
therefore distinct. Their values are in $[0,1]$, so they solve both the
polynomial Fisher--KPP equation and \eqref{eq:problem} with the global
reaction \eqref{eq:reaction}. This proves Theorem~\ref{thm:main}.

\paragraph{AI use.}
OpenAI Codex assisted with the mathematical development, source checking and
preparation of the manuscript. Two further fresh Codex contexts supplied
reader and referee feedback on an earlier draft, followed by an author
revision. These were AI reviews; independent mathematical review is requested.
The author is responsible for the contents.

\begin{thebibliography}{9}
\bibitem{BGKPP} H. Berestycki and C. Graham,
\emph{The steady states of strong-KPP reactions in general domains},
arXiv:2212.06611 (2022), \url{https://arxiv.org/abs/2212.06611};
Journal of the European Mathematical Society, published online 2025,
\href{https://doi.org/10.4171/JEMS/1578}{doi:10.4171/JEMS/1578}.
\bibitem{BGSC} H. Berestycki and C. Graham,
\emph{A Stable-Compact Method for Qualitative Properties of Semilinear Elliptic Equations},
Archive for Rational Mechanics and Analysis \textbf{250} (2026), article 21,
\href{https://doi.org/10.1007/s00205-026-02168-6}{doi:10.1007/s00205-026-02168-6}.
\bibitem{RT} L. Rossi and O. Tough,
\emph{Sharp dimension threshold for unique FKPP steady states and BBM recurrence},
arXiv:2610.17121v1 (2026), \url{https://arxiv.org/abs/2610.17121v1}.
\bibitem{GT} D. Gilbarg and N. S. Trudinger,
\emph{Elliptic Partial Differential Equations of Second Order},
Classics in Mathematics, Springer, 2001,
\href{https://doi.org/10.1007/978-3-642-61798-0}{doi:10.1007/978-3-642-61798-0}.
\end{thebibliography}
\end{document}
